% MeLLM @ ACL 2026 — blind review submission
% Compile:  pdflatex main.tex
%           bibtex main
%           pdflatex main.tex  (twice for cross-refs)
%
% ─────────────────────────────────────────────────────────────────────────────
\documentclass[11pt]{article}
\usepackage{acl}
\usepackage{times}

% Encoding (pdfLaTeX — no fontspec/polyglossia needed)
% T2A loaded alongside T1 to enable Cyrillic in Table A1 via \ru{}; T1 remains default
\usepackage[T2A,T1]{fontenc}
\usepackage[utf8]{inputenc}
% Cyrillic support for Table A1 (switches encoding locally, does not affect body text)
\newcommand{\ru}[1]{{\fontencoding{T2A}\fontfamily{cmr}\selectfont #1}}
\usepackage{microtype}

% Math
\usepackage{amsmath,amssymb}

% Tables
\usepackage{booktabs}
\usepackage{multirow}

% Figures
\usepackage{graphicx}
\usepackage{xcolor}

% Compact lists
\usepackage{enumitem}

% Fix: suppress lineno line numbers inside float environments (review mode artifact)
\usepackage{etoolbox}
\BeforeBeginEnvironment{table}{\begin{nolinenumbers}}
\AfterEndEnvironment{table}{\end{nolinenumbers}}
\BeforeBeginEnvironment{table*}{\begin{nolinenumbers}}
\AfterEndEnvironment{table*}{\end{nolinenumbers}}
\BeforeBeginEnvironment{figure}{\begin{nolinenumbers}}
\AfterEndEnvironment{figure}{\end{nolinenumbers}}
\BeforeBeginEnvironment{figure*}{\begin{nolinenumbers}}
\AfterEndEnvironment{figure*}{\end{nolinenumbers}}

% Graceful figure inclusion: shows a grey placeholder if the PDF is absent
\newcommand{\safefig}[3][width=\linewidth]{%
  \IfFileExists{#2}%
    {\includegraphics[#1]{#2}}%
    {\fbox{\begin{minipage}{0.96\linewidth}\centering\vspace{0.7cm}%
      \textcolor{gray}{\small[Figure: #3]}\vspace{0.7cm}%
    \end{minipage}}}%
}

% ─────────────────────────────────────────────────────────────────────────────
\title{Causal Localization of the English Pivot in LLaVA:\\
Mechanistic VLM Analysis and Training-Free\\
Multilingual Steering\thanks{Code: \url{https://github.com/azraihan/mul-vlm-mech-interp}}}

\author{
  \textbf{Abrar Zahin Raihan}$^*$ \and \textbf{Aurchi Chowdhury}$^*$ \\
  Bangladesh University of Engineering and Technology \\
  \texttt{raihanzahin95@gmail.com} \quad \texttt{aurchichowdhury07@gmail.com} \\[2pt]
  {\small $^*$Equal contribution}
}

\begin{document}
\maketitle

% ─────────────────────────────────────────────────────────────────────────────
\begin{abstract}
Multilingual vision-language models (VLMs) consistently underperform on
non-English visual queries, yet the internal mechanism behind this disparity
remains unknown.
As a focused case study on LLaVA-1.5-7B, we apply logit-lens analysis and
causal activation patching to show that non-English visual queries are routed
through an English-biased representational bottleneck in layers 5--17,
extending the English-pivot phenomenon of \citet{wendler-etal-2024-llamas} to
the multimodal setting.
Peak causal influence occurs at layer~8 ($\overline{\mathrm{AIE}}{=}0.49$,
averaged across languages), with all measurable pivot signal running through
text-token positions.
Without meaningful visual content (blank-image condition), language-specific
representations do not emerge at any layer, showing that the pivot is
image-content-dependent rather than triggered by any visual input.
Building on these findings, we derive training-free language-steering vectors
at the mechanistically identified pivot layers, improving Russian VQA by
+6.5 percentage points (pp) and Portuguese by +4.0~pp on MMMB without any
fine-tuning---the latter surpassing the English baseline.
Within this case study, our results are consistent with the English pivot
being a structural property of the LLM backbone that multimodal pre-training
does not mitigate; extending this mechanistic methodology to other VLMs and
language families remains an important direction for future work.
\end{abstract}

% ═══════════════════════════════════════════════════════════════════════════════
\section{Introduction}

Multilingual VLMs are deployed globally, yet a persistent gap separates their
English and non-English performance on visual question answering
\citep{fan-etal-2025-plast,song-etal-2025-missing}.
Prior work attributes this gap to training-data imbalance or insufficient
multilingual instruction tuning, and proposes remedies such as language-aware
fine-tuning \citep{fan-etal-2025-plast} or culturally diverse data
\citep{nyandwi-etal-2025-grounding}.
Yet no prior work has asked the mechanistic question: \emph{where} in the
network does English bias arise, and what causal role does the visual modality
play?

The text-only interpretability literature provides a starting hypothesis.
\citet{wendler-etal-2024-llamas} show that LLaMA-family models process
multilingual input through three phases: an \emph{input space}, a
\emph{concept space} where representations become English-biased, and an
\emph{output space} where the target language is reconstructed.
\citet{ferrando-costa-jussa-2024-similarity} extend this to circuit-level
analysis across languages.
Neither work addresses VLMs, leaving open whether the visual modality interacts
with or suppresses this pivot.

We fill this gap with the first causal mechanistic analysis of the English
pivot in a multilingual VLM, presenting a focused case study on LLaVA-1.5-7B.
We treat our findings as specific to this model and architecture; whether
analogous structures arise in other VLMs is an open question we leave to
future work.
Our contributions are:

\begin{itemize}[noitemsep,topsep=2pt]
\item We confirm that LLaVA-1.5-7B routes non-English visual queries through
  an English-biased concept space in \textbf{layers 5--17}, with peak causal
  responsibility at layer~8 (cross-language average $\overline{\mathrm{AIE}}{=}0.49$; \S\ref{sec:patching}).
\item We show that the pivot is image-content-dependent: language-specific
  representations do not emerge without semantically rich visual input
  (blank-image condition; \S\ref{sec:logit_lens}).
\item Under our patching design, all measurable pivot signal runs through
  text-token positions; isolating visual causal contributions would require
  paired semantically-equivalent cross-lingual images (\S\ref{sec:patching}).
\item We derive training-free steering vectors at the mechanistically
  identified pivot layers, yielding \textbf{+6.5~pp} on Russian and
  \textbf{+4.0~pp} on Portuguese on the MMMB multiple-choice benchmark,
  without any fine-tuning (\S\ref{sec:steering}).
\end{itemize}

% ═══════════════════════════════════════════════════════════════════════════════
\section{Related Work}

\paragraph{Mechanistic interpretability of multilingual LLMs.}
\citet{wendler-etal-2024-llamas} establish the three-phase English-pivot model
for LLaMA via logit-lens analysis.
\citet{dumas-etal-2024-llamas} use activation patching to show that concept
representations are language-agnostic at intermediate layers, separating the
underlying concept from its surface language form.
\citet{ferrando-costa-jussa-2024-similarity} show language-agnostic circuits
are reused across languages, implying the pivot is structural.
\citet{brinkmann-etal-2025-large} use causal sparse-autoencoder feature analysis
to show shared latent grammatical representations across typologically diverse
languages; \citet{resck-etal-2025-explainability} survey mechanistic and
interpretability findings across multilingual LLMs more broadly.
All of these works study text-only models; we extend causal mechanistic
analysis to VLMs.

\paragraph{Multilingual vision-language models.}
Multilingual VLM architectures such as mBLIP \citep{geigle-etal-2023-mblip}
and PALO \citep{maaz-etal-2024-palo} extend English-centric models through
multilingual encoders or instruction tuning, but do not examine internal mechanisms.
\citet{fan-etal-2025-plast} propose PLAST, a language-aware approach requiring
fine-tuning.
\citet{song-etal-2025-missing} document non-English VQA failure modes without
investigating internal mechanisms.
\citet{nyandwi-etal-2025-grounding} address cultural diversity at the data
level.
In contrast, we provide a causal mechanistic account and a training-free fix.

\paragraph{Mechanistic interpretability for VLMs.}
\citet{golovanevsky-etal-2025-notice} apply activation patching to LLaVA and
BLIP to study how visual information is integrated under image corruption,
finding that LLaVA's attention heads perform outlier suppression rather than
visual grounding---a complementary analysis of the visual integration pathway
rather than the language-processing pathway we study.
Our steering intervention is related to representation engineering
\citep{zou-etal-2023-representation}, but derived from mechanistically
identified pivot layers rather than contrast pairs.

% ═══════════════════════════════════════════════════════════════════════════════
\section{Background}
\label{sec:background}

\paragraph{LLaVA-1.5-7B.}
LLaVA-1.5-7B \citep{liu-etal-2024-improved} connects a CLIP ViT-L/14 visual
encoder to a Vicuna-7B backbone via a two-layer MLP projector.
An image is encoded into 576 patch tokens (positions 1--576 in the sequence),
followed by the tokenised question.
The backbone has 32 Llama-2 decoder layers of hidden size 4096.

\paragraph{Logit lens and PMI.}
The logit lens \citep{nostalgebraist-2020} projects an intermediate hidden
state $\mathbf{h}_\ell$ through the final layer norm and LM head to obtain
token probabilities at any depth.
We report \emph{pointwise mutual information}:
$\mathrm{PMI}_\ell[\tau] = \log p_\ell[\tau] - \log p_0[\tau]$,
where $p_0$ is the LM-head output at a zero hidden state (the model's
unconditional prior), to remove vocabulary-frequency bias.

\paragraph{Causal tracing and AIE.}
Causal tracing \citep{meng-etal-2022-locating} replaces the clean hidden
state at layer $\ell$ with the corresponding state from a corrupted run,
measuring the output change.
We define the \emph{Average Indirect Effect}:
\begin{equation}
  \mathrm{AIE}(\ell)
  = \mathbb{E}_e\!\left[
      p_\ell^{\text{patch}}\!\left[\tau_{\mathrm{en}}\right]
    - p^{\text{clean}}\!\left[\tau_{\mathrm{en}}\right]
    \right],
\end{equation}
where the clean run uses a non-English question and the corrupted run uses the
English equivalent on the same image.
$\tau_{\mathrm{en}}$ is the first token of the model's English response,
determined dynamically per example (typically ``The'', id~450; see
Appendix~\ref{app:patching}).
Positive AIE at layer $\ell$ means that layer causally carries English-pivot
information.

% ═══════════════════════════════════════════════════════════════════════════════
\section{Experiments and Results}
\label{sec:experiments}

\subsection{Experimental Setup}
\label{sec:setup}

\paragraph{Model.}
We use \texttt{llava-hf/llava-1.5-7b-hf} in float16, no quantisation or
fine-tuning, greedy decoding.

\paragraph{Probing set.}
We build a 195-example probing set from COCO val2017
\citep{lin-etal-2014-microsoft} across 15 object categories (person, car,
bus, train, boat, orange, book, clock, knife, spoon, umbrella, bird, cat,
dog, apple).
For each example we ask the question in four languages: English (en),
Portuguese (pt), German (de), and Russian (ru), using forced-choice prompts
listing all 15 category names in the target language.
We track the first sub-token of each answer word using contextual-subtraction
tokenisation to match generation-time tokenisation (Appendix~\ref{app:probing}).
This controlled object-naming paradigm is designed to identify the earliest
layers at which the model develops language-specific representations under a
semantically invariant visual stimulus; the transfer of vectors derived from
this set to the broader MMMB benchmark is an empirical finding.

\paragraph{Evaluation benchmarks.}
We evaluate steering on MMMB (Massive Multilingual Multimodal Benchmark;
\citealp{sun-etal-2023-mmmb}), a multilingual multiple-choice VQA benchmark
covering English, Portuguese, and Russian
(200 examples each).
German is included in the mechanistic analysis (Figures~\ref{fig:logit_lens}--\ref{fig:visual})
and cross-lingual transfer ablation (\S\ref{sec:ablations}) because its
residual-stream AIE profile peaks at layer~10 with a pivot region of
approximately 5--17, closely matching the range identified for all languages---making
it a natural cross-lingual transfer control.
MMMB does not provide a German split, so German steering is assessed indirectly via transfer.

\paragraph{Baselines.}
We compare our steered model ($\alpha{>}0$) against unmodified LLaVA-1.5-7B
(baseline, $\alpha{=}0$) and against steering vectors extracted from fixed
layer ranges---early (0--10), mid (11--21), and late (22--31)---to test
whether mechanistic localisation is necessary.

% ─────────────────────────────────────────────────────────────────────────────
\subsection{The English Pivot in VLMs}
\label{sec:logit_lens}

\begin{figure*}[t]
  \centering
  \safefig[width=0.88\textwidth]{../outputs/figures/fig1_logit_lens.pdf}%
         {Logit Lens PMI across 32 layers × 3 languages × 2 conditions}
  \caption{
    \textbf{Logit-lens PMI across layers.}
    Blue (solid): English answer token; red (dashed): target-language token.
    Top row: real image; bottom row: blank image.
    Grey band: mechanistically identified pivot region (layers 5--17).
    With a real image, both tokens develop strong positive PMI in late layers
    (20--29), with English peaking first at layers 20--23.
    Without meaningful visual content, PMI is predominantly negative at every
    layer; where marginally positive, values are substantially weaker than in
    the image-present condition---language-specific representations do not emerge.
  }
  \label{fig:logit_lens}
\end{figure*}

Figure~\ref{fig:logit_lens} shows logit-lens PMI across all 32 layers.
German is included in all mechanistic figures as a cross-lingual control;
it has no MMMB split, so its steering effect is assessed indirectly via
cross-lingual transfer (\S\ref{sec:ablations}).

\paragraph{Image-present condition.}
With a real image, both English and target-language PMI grow substantially in
late layers.
English PMI peaks at layers 20--23 (values 8.1--8.4) across all three
languages, consistently \emph{before} the target-language peak at layers 27--29
(values 8.6--10.1).
This ordering---English representations crystallising 5--7 layers before the
target language---directly instantiates Wendler et al.'s concept-space phase
in the multimodal setting.
The grey pivot region (layers 5--17) corresponds to where English PMI first
begins to build, preceding both peaks.

\paragraph{Blank-image condition.}
Without meaningful visual content (image replaced by a blank tensor of
the same spatial dimensions), PMI values are predominantly negative at every
layer for every language; where marginally positive, values are substantially
weaker than in the image-present condition.
This shows that the pivot is image-content-dependent: it is not triggered by
any visual input, but requires semantically rich visual content for
language-specific representations to form.

\noindent\textbf{Logit-lens vs.\ causal patching.}\enspace
These measurements are complementary: the logit lens shows where representations
\emph{crystallise} into decodable probabilities (a readout property), while
causal patching (\S\ref{sec:patching}) shows where pivot information
\emph{transfers} causally---hence the pivot region (5--17) precedes the
late-layer PMI peaks (20--29).

% ─────────────────────────────────────────────────────────────────────────────
\subsection{Causal Localisation via Activation Patching}
\label{sec:patching}

\begin{figure*}[t]
  \centering
  \safefig[width=\textwidth]{../outputs/figures/fig2_patching_heatmap.pdf}%
         {AIE heatmap: residual / attn / MLP}
  \caption{
    \textbf{Average Indirect Effect (AIE) heatmap.}
    Rows: Portuguese, German, Russian.
    Columns: residual stream / attention output / MLP output.
    The residual-stream pivot at layers 5--17 is pronounced across all
    languages; attention and MLP patches carry negligible weight.
    Yellow dashed box: shared pivot region (layers 5--17) used for steering.
    Green dashed box: language-specific pivot region inferred from each
    language's AIE profile, shown to highlight the per-language variation
    (e.g.\ Portuguese tapers off around layer~14 while Russian extends to~17).
  }
  \label{fig:patching_heatmap}
\end{figure*}

\begin{figure}[t]
  \centering
  \safefig[width=0.85\linewidth]{../outputs/figures/fig3_visual_contribution.pdf}%
         {Full vs. visual-token-only AIE}
  \caption{
    \textbf{Design validation: visual-token-only AIE is identically zero.}
    The same image appears in both passes, so the 576 CLIP patch activations
    are identical by construction and contribute zero differential signal.
    This confirms correct implementation: all measurable pivot signal under
    this design runs through text-token positions.
  }
  \label{fig:visual}
\end{figure}

\paragraph{Residual stream patching.}
Figure~\ref{fig:patching_heatmap} (left column) shows strongly positive AIE
in layers 5--17 and near-zero thereafter.
Peak cross-language average AIE is at \textbf{layer~8}
($\overline{\mathrm{AIE}}{=}0.49$); per-language peaks reach 0.52 (PT,
layer~8), 0.62 (DE, layer~10), and 0.85 (RU, layer~10), reflecting
Russian's substantially broader and stronger pivot.
The signal remains above the 50\%-of-peak threshold (0.24) through layer~17
and drops to noise floor ($<0.04$) by layer~20.
We define the \textbf{pivot region} as layers 5--17.

The pivot region occupies the first half of the LLM, consistent with
Wendler et al.'s concept-space phase, and precedes the late-layer (20--29)
representational peaks observed in \S\ref{sec:logit_lens}.

\paragraph{Attention and MLP patches.}
Middle and right columns of Figure~\ref{fig:patching_heatmap} show near-zero
AIE for both attention-output and MLP-output patches across all layers
(max $|\overline{\mathrm{AIE}}|{<}0.01$ for attention; ${<}0.01$ for MLP
except layer~0 where MLP${\approx}0.08$, reflecting sub-word tokenisation
effects \citep{geva-etal-2022-transformer}).
The English pivot operates through the integrated residual stream rather than
through separable attention heads or feed-forward networks.

\paragraph{Design validation: visual token isolation.}
Figure~\ref{fig:visual} serves as a design validation rather than a mechanistic
finding.
Because the same image is used in both the non-English (clean) and English
(corrupted) runs, the 576 CLIP visual-token activations are identical by
construction and contribute zero differential signal.
The identically-zero visual-token curve confirms that our patching setup
correctly isolates text-token contributions---any non-zero value would indicate
a confound in the implementation.
What the design \emph{does} establish is that the pivot's causal pathway does
not depend on any language-specific visual signal; the blank-image ablation
(\S\ref{sec:logit_lens}) provides complementary evidence that the pivot is
image-content-dependent---semantically rich visual input is required for
language-specific representations to form.
Whether visual tokens could independently carry pivot information under a
paired-image design (e.g.\ semantically equivalent images in different
languages) remains an open question for future work.

% ─────────────────────────────────────────────────────────────────────────────
\subsection{Training-Free Steering at Pivot Layers}
\label{sec:steering}

\paragraph{Steering vector extraction.}
For each non-English language $\ell$, we extract a steering vector at each
pivot layer $L \in \{5,\ldots,17\}$:
\begin{equation}
  \mathbf{v}_L^{(\ell)}
  = \frac{1}{N}\sum_{i=1}^{N}
    \mathbf{h}_L^{(\ell,i)} - \mathbf{h}_L^{(\mathrm{en},i)},
\end{equation}
where $\mathbf{h}_L^{(\cdot,i)}$ is the residual-stream hidden state at the
last answer position for example $i$ ($N{=}195$).
Split-half cosine similarity exceeds 0.85 across all languages and pivot
layers, indicating stable steering directions.

\paragraph{Steered inference.}
At inference time we add $\alpha\,\mathbf{v}_L^{(\ell)}$ to the hidden state
at each pivot layer via an in-place forward hook, applied to the last active
position at every decode step (KV-cache compatible;
Appendix~\ref{app:steering}).

\begin{table}[t]
\centering
\small
\setlength{\tabcolsep}{5pt}
\caption{
  \textbf{Main results on MMMB} (4-way multiple-choice accuracy, \%).
  ``Ours'': best accuracy over
  $\alpha \in \{0.01, 0.02, 0.05, 0.07, 0.1, 0.2, 0.3, 0.4, 0.5\}$
  (optimum $\alpha{=}0.2$ for PT and RU independently).
  EN: baseline only (no steering vector for English).
  $\Delta$: gain over baseline.
}
\label{tab:main}
\begin{tabular}{lccc}
\toprule
\textbf{Method} & \textbf{EN} & \textbf{PT} & \textbf{RU} \\
\midrule
Baseline              & 66.0 & 63.0 & 51.5 \\
Ours ($\alpha{=}0.2$) & 66.0 & \textbf{67.0} & \textbf{58.0} \\
\midrule
$\Delta$              & ---  & +4.0 & +6.5 \\
\bottomrule
\end{tabular}
\end{table}

\paragraph{Results (Table~\ref{tab:main}).}
Steering yields \textbf{+6.5~pp} on Russian MMMB (51.5\% $\to$ 58.0\%
at $\alpha{=}0.2$), substantially reducing the 14.5-point gap between Russian and
English baselines.
Portuguese gains \textbf{+4.0~pp} (63.0\% $\to$ 67.0\%), surpassing the English
baseline of 66.0\%---demonstrating that mechanistically guided steering can lift
non-English performance above the monolingual baseline without any fine-tuning.
The Portuguese AIE profile peaks earlier and tapers off around layer~14
(Figure~\ref{fig:patching_heatmap}), suggesting that language-specific pivot
ranges (e.g.\ 5--14 for Portuguese) may yield further gains; we leave this
optimisation as a direction for future work.
Figure~\ref{fig:case_study} shows a representative Portuguese example where
steering shifts the model's prediction from the wrong option (C) to the
correct answer (D), with the PMI curves at the pivot layers revealing the
mechanism behind the correction.

\begin{figure*}[t]
  \centering
  \safefig[width=0.88\textwidth]{../outputs/figures/pt_0015.png}%
         {Portuguese MMMB case study: baseline vs.\ steered PMI curves}
  \caption{
    \textbf{Representative Portuguese MMMB case (steered wins).}
    \emph{Left:} input image and question with four answer options; correct
    answer is D.
    The question asks about the colour of a bathroom fixture---a
    colour-attribute question, not an object-naming question, demonstrating
    that steering vectors transfer to qualitatively different question types
    beyond the probing-set distribution.
    \emph{Top right (Baseline):} logit-lens PMI curves for all four options
    across 32 layers---the model incorrectly peaks on option~C at layer~31.
    \emph{Bottom right (Steered, $\alpha{=}0.2$):} after applying the
    Portuguese steering vector at pivot layers 5--17, option~D dominates
    at layer~31, recovering the correct answer.
    The shaded fill in the pivot region highlights where C and D compete;
    steering resolves this competition in favour of the correct option.
  }
  \label{fig:case_study}
\end{figure*}

% ─────────────────────────────────────────────────────────────────────────────
\subsection{Ablations}
\label{sec:ablations}

\begin{table}[t]
\centering
\small
\caption{
  \textbf{Layer range ablation on MMMB (best $\alpha$ per cell, \%).}
  Best $\alpha$ shown in parentheses where it differs from 0.2.
}
\label{tab:layer_ablation}
\begin{tabular}{lcc}
\toprule
\textbf{Layer range} & \textbf{PT} & \textbf{RU} \\
\midrule
Early (0--10)    & 63.5 ($\alpha{=}0.2$)  & 51.5 ($\alpha{=}0.01$) \\
Mid   (11--21)   & 63.0 ($\alpha{=}0.01$) & 54.5 ($\alpha{=}0.1$)  \\
Late  (22--31)   & 63.0 ($\alpha{=}0.01$) & 51.5 ($\alpha{=}0.01$) \\
Ours  (5--17)    & \textbf{67.0} ($\alpha{=}0.2$) & \textbf{58.0} ($\alpha{=}0.2$) \\
\bottomrule
\end{tabular}
\end{table}

\begin{table}[t]
\centering
\small
\caption{
  \textbf{Visual token ablation on MMMB (\%).}
  ``Steer'' = best $\alpha$: $\alpha{=}0.2$ (real), $\alpha{=}0.02$ (blank).
}
\label{tab:visual_ablation}
\resizebox{\linewidth}{!}{%
\begin{tabular}{lcccc}
\toprule
\textbf{Condition} & \textbf{PT\,Base} & \textbf{PT\,Steer}
                   & \textbf{RU\,Base} & \textbf{RU\,Steer} \\
\midrule
Real image  & 63.0 & 67.0 & 51.5 & 58.0 \\
Blank image & 42.0 & 41.5 & 46.5 & 48.5 \\
\bottomrule
\end{tabular}}
\end{table}

\begin{table}[t]
\centering
\small
\caption{
  \textbf{Cross-lingual transfer on MMMB (best $\alpha$ per cell, \%).}
}
\label{tab:cross_lingual}
\begin{tabular}{lcc}
\toprule
\textbf{Vector source} & \textbf{MMMB-PT} & \textbf{MMMB-RU} \\
\midrule
PT (own) & ---                   & 55.5 ($\alpha{=}0.3$) \\
DE       & 67.0 ($\alpha{=}0.2$) & 58.0 ($\alpha{=}0.2$) \\
RU (own) & 63.5 ($\alpha{=}0.3$) & ---  \\
\bottomrule
\end{tabular}
\end{table}

\begin{figure}[t]
  \centering
  \safefig[width=\linewidth]{../outputs/figures/fig4_alpha_curves.pdf}%
         {Alpha sensitivity curves (MMMB)}
  \caption{
    \textbf{Sensitivity to $\alpha$ on MMMB.}
    Performance peaks at $\alpha{=}0.2$ and degrades sharply beyond
    $\alpha{=}0.3$, consistent with prior representation-engineering results
    \citep{zou-etal-2023-representation}.
  }
  \label{fig:alpha}
\end{figure}

\begin{figure*}[t]
  \centering
  \safefig[width=0.55\textwidth]{../outputs/figures/fig5_qualitative.pdf}%
         {Qualitative logit-lens walkthrough with and without steering}
  \caption{
    \textbf{Qualitative logit-lens walkthrough (probing set).}
    Russian (top), German (middle), Portuguese (bottom) without steering (left)
    and with $\alpha{=}0.2$ (right).
    Red arrow: layer where target-language PMI first exceeds English.
    Steering shifts the crossover earlier, amplifying target-language
    representations at pivot layers.
  }
  \label{fig:qualitative}
\end{figure*}

\paragraph{Layer range ablation (Table~\ref{tab:layer_ablation}).}
Our mechanistic pivot region (5--17) outperforms all fixed ranges on both
languages: +4.0~pp vs.\ +0.5~pp (early) for Portuguese, and +6.5~pp vs.\
+3.0~pp (mid) for Russian.
Early and late ranges match the baseline for Russian; mid provides +3.0~pp
but falls well short of our range's +6.5~pp.
This ablation shows that mechanistic localisation beats arbitrary partitioning;
we note that any AIE-informed selection of similar width would likely achieve
comparable gains, and the key value of the mechanistic analysis is identifying
the pivot region rather than optimising layer boundaries post-hoc.

\paragraph{Visual token ablation (Table~\ref{tab:visual_ablation}).}
Blank images drastically lower baseline performance (PT: $-21$~pp;
RU: $-5$~pp), consistent with the image-content-dependence finding from
\S\ref{sec:logit_lens}.
With blank images, steering is marginal or slightly harmful, as language-specific
representations require semantically rich visual content to form.

\paragraph{Steering coefficient sensitivity (Figure~\ref{fig:alpha}).}
Performance peaks at $\alpha{=}0.2$ for both MMMB languages (independently)
and degrades sharply beyond $\alpha{=}0.3$.
Large $\alpha$ values corrupt generation quality (MMMB-RU drops to 0\% at
$\alpha{=}0.5$, indicating degenerate outputs).

\paragraph{Cross-lingual transfer (Table~\ref{tab:cross_lingual}).}
The \textbf{German steering vector transfers perfectly} to both Portuguese
(67.0\% = own-language) and Russian (58.0\% = own-language).
In contrast, the Portuguese vector applied to Russian yields 55.5\%
($-2.5$~pp vs.\ own-language).
The full transferability of the German vector suggests it captures a general
``non-English'' axis rather than a language-specific direction, consistent
with the shared concept space observed in activation patching.

% ═══════════════════════════════════════════════════════════════════════════════
\section{Discussion}

\paragraph{Why does the pivot exist in VLMs?}
The pivot is a backbone property: Vicuna's English-dominant pretraining maps
non-English sequences into an English-centric concept space.
Because CLIP is language-agnostic, the MLP projector translates visual features
into the backbone's embedding space without any language-direction component---visual
tokens activate language-specific processing but cannot signal \emph{which}
language to use, explaining the zero visual-token AIE.
Multimodal training optimises only the projector and leaves the backbone's
language topology intact.

\paragraph{Steering direction geometry.}
The complete transferability of the German vector suggests that pt, de, and ru
steering directions approximately coincide---a single ``non-English'' axis
rather than three language-specific directions---consistent with the shared
concept-space hypothesis of \citet{wendler-etal-2024-llamas} and raising the
possibility that one universal vector could suffice for multilingual correction.

\paragraph{Practical Implications.}
Our steering approach is KV-cache compatible and requires no weight updates,
making it a drop-in patch for existing LLaVA-1.5 deployments serving
non-English users on structured VQA tasks.
The complete cross-lingual transfer of the German vector suggests that
a single ``non-English'' vector---computed once across languages---may
generalise to unseen non-English inputs, potentially reducing the
per-language probing requirement; we leave this as an empirical question
for future work.
More broadly, the logit-lens and AIE pipeline could serve as a lightweight
pre-deployment audit for new VLMs, localising whether an English pivot
exists and at which layers it peaks, complementing rather than replacing
full multilingual benchmarking.

\paragraph{Limitations.}
This is a single-model case study; generalisation to other VLMs and language
families remains open.
The steering coefficient $\alpha$ is selected on the MMMB test set, so reported
gains are upper bounds on out-of-sample performance.
The shared pivot region (5--17) may be over-inclusive for some languages
(Portuguese AIE tapers at layer~14); per-language ranges are a promising
direction for stronger gains.
Finally, the zero visual-token AIE is a design artefact of using the same
image in both passes; paired cross-lingual images would probe whether visual
tokens can independently carry pivot information.

% ═══════════════════════════════════════════════════════════════════════════════
\section{Conclusion}

LLaVA-1.5-7B routes non-English visual queries through an English-biased
representational bottleneck in layers 5--17 (peak $\overline{\mathrm{AIE}}{=}0.49$
at layer~8), inherited from the Vicuna backbone and unmitigated by multimodal
pre-training.
The pivot is image-content-dependent---language-specific representations
require semantically rich visual content to emerge---but is mediated
entirely through text-token positions.
Training-free steering at the mechanistically identified layers yields
+6.5~pp on Russian and +4.0~pp on Portuguese (surpassing English), with
German vectors transferring perfectly---establishing the pivot as structural
and providing a blueprint for mechanistic analysis of other VLMs.

% ─────────────────────────────────────────────────────────────────────────────
\bibliography{references}

% ═══════════════════════════════════════════════════════════════════════════════
\appendix

\section{Probing Set Details}
\label{app:probing}

Table~\ref{tab:probing_categories} lists the 15 COCO categories and their
Portuguese, German, and Russian answer tokens.
All answer words are tracked by their first sub-token under the Vicuna
tokeniser, using contextual subtraction (encoding ``ASSISTANT: \{word\}'' and
stripping the prefix tokens) to match generation-time tokenisation.
The leading-space token (id~29871) is filtered before selecting the first
sub-token.

\begin{table}[h]
\centering
\footnotesize
\setlength{\tabcolsep}{4pt}
\caption{
  Object categories with PT, DE, and RU answer tokens.
  Multi-token words tracked by first sub-token only ($\dagger$).
  All Russian (Cyrillic) entries are multi-token under the Vicuna tokeniser.
}
\label{tab:probing_categories}
\begin{tabular}{llll}
\toprule
\textbf{Category} & \textbf{PT} & \textbf{DE} & \textbf{RU} \\
\midrule
person   & pessoa              & Person              & \ru{человек}$^\dagger$ \\
car      & carro               & Auto                & \ru{машина}$^\dagger$ \\
bus      & \^{o}nibus$^\dagger$ & Bus                & \ru{автобус}$^\dagger$ \\
train    & trem                & Zug                 & \ru{поезд}$^\dagger$ \\
boat     & barco               & Boot                & \ru{лодка}$^\dagger$ \\
orange   & laranja$^\dagger$   & Orange              & \ru{апельсин}$^\dagger$ \\
book     & livro               & Buch                & \ru{книга}$^\dagger$ \\
clock    & rel\'ogio$^\dagger$ & Uhr                 & \ru{часы}$^\dagger$ \\
knife    & faca                & Messer$^\dagger$    & \ru{нож}$^\dagger$ \\
spoon    & colher$^\dagger$    & L\"offel$^\dagger$  & \ru{ложка}$^\dagger$ \\
umbrella & guarda-chuva$^\dagger$ & Regenschirm$^\dagger$ & \ru{зонт}$^\dagger$ \\
bird     & p\'assaro$^\dagger$ & Vogel$^\dagger$     & \ru{птица}$^\dagger$ \\
cat      & gato                & Katze$^\dagger$     & \ru{кот}$^\dagger$ \\
dog      & cachorro$^\dagger$  & Hund                & \ru{собака}$^\dagger$ \\
apple    & ma\c{c}\~a$^\dagger$ & Apfel$^\dagger$   & \ru{яблоко}$^\dagger$ \\
\bottomrule
\end{tabular}
\end{table}

\section{Implementation Details}
\label{app:patching}

Code will be made available upon acceptance.

\paragraph{Sequence alignment.}
Non-English questions are 2--5 tokens longer than English counterparts.
We keep the clean (non-English) input at full length and truncate only the
corrupted (English) input to \texttt{min\_len} to preserve the correct
\texttt{answer\_pos}.

\paragraph{Pre-hook patching under transformers 5.0.}
\texttt{GradientCheckpointingLayer} wrappers in transformers 5.0 discard
post-hook return values, so we register patching as a \emph{pre-hook} on
layer $\ell{+}1$ rather than a post-hook on layer $\ell$.
For layer $\ell{=}31$, the pre-hook is on
\texttt{model.language\_model.model.norm}.

\paragraph{Response-token target and steering.}
\label{app:steering}
$\tau_\mathrm{en}$ is determined dynamically via one-token greedy decoding on
the English input (consistently ``The'', id~450).
Steering applies $\mathbf{h}_L \mathrel{+}= \alpha \cdot \mathbf{v}_L^{(\ell)}$
at the last active token position at each decode step via a forward hook on
\texttt{model.language\_model.model.layers[L]}, covering both prefill and
autoregressive steps under KV-cache.

\end{document}